\section{Lat-Lon Dataset Groupings}
\subsection{latlon dataset SEA\_SURFACE\_HEIGHT}
\newp
\subsubsection{Overview}
This dataset provides 2D fields of dynamic sea surface height interpolated to a regular 0.5-degree grid from the ECCO Version 4 Release 4 (V4r4) ocean and sea-ice state estimate. The dataset is provided on daily-average and monthly-average time resolution.
\\\\
Note: SSH (dynamic sea surface height) = SSHNOIBC (dynamic sea surface without the inverse barometer correction) - SSHIBC (inverse barometer correction). The inverted barometer correction accounts for variations in sea surface height due to atmospheric pressure variations.
\\
\begin{longtable}{|m{0.15\textwidth}|m{0.64\textwidth}|m{0.12\textwidth}|}
\caption{Coordinates and Variables in the dataset SEA\_SURFACE\_HEIGHT}
\label{tab:table-<xarray.Dataset> Size: 3MB
Dimensions:         (time: 1, latitude: 360, longitude: 720, nv: 2)
Coordinates:
  * time            (time) datetime64[ns] 8B 2017-12-29T12:00:00
  * latitude        (latitude) float32 1kB -89.75 -89.25 -88.75 ... 89.25 89.75
  * longitude       (longitude) float32 3kB -179.8 -179.2 -178.8 ... 179.2 179.8
    time_bnds       (time, nv) datetime64[ns] 16B ...
    latitude_bnds   (latitude, nv) float32 3kB ...
    longitude_bnds  (longitude, nv) float32 6kB ...
Dimensions without coordinates: nv
Data variables:
    SSH             (time, latitude, longitude) float32 1MB ...
    SSHIBC          (time, latitude, longitude) float32 1MB ...
    SSHNOIBC        (time, latitude, longitude) float32 1MB ...
Attributes: (12/57)
    acknowledgement:              This research was carried out by the Jet Pr...
    author:                       Ian Fenty and Ou Wang
    cdm_data_type:                Grid
    comment:                      Fields provided on a regular lat-lon grid. ...
    Conventions:                  CF-1.8, ACDD-1.3
    coordinates_comment:          Note: the global 'coordinates' attribute de...
    ...                           ...
    time_coverage_duration:       P1D
    time_coverage_end:            2017-12-30T00:00:00
    time_coverage_resolution:     P1D
    time_coverage_start:          2017-12-29T00:00:00
    title:                        ECCO Sea Surface Height - Daily Mean 0.5 De...
    uuid:                         7beb77e0-404a-11eb-8bf7-0cc47a3f891b-fields} \\ 
\hline \endhead \hline \endfoot
\rowcolor{lightgray} \multicolumn{1}{|c|}{\textbf{Coordinates}} & \multicolumn{1}{|c|}{\textbf{Description of data coordinates}} &  \multicolumn{1}{|c|}{\textbf{Unit}}\\ \hline
time &Center time of averaging period &--none--  \\ \hline
latitude &Latitude at grid cell center &degrees\_north  \\ \hline
longitude &Longitude at grid cell center &degrees\_east  \\ \hline
time\_bnds &Time bounds of averaging period &--none--  \\ \hline
latitude\_bnds &Latitude bounds grid cells &--none--  \\ \hline
longitude\_bnds &Longitude bounds grid cells &--none--  \\ \hline
\rowcolor{lightgray} \multicolumn{1}{|c|}{\textbf{Variables}} & \multicolumn{1}{|c|}{\textbf{Description of data variables}} &  \multicolumn{1}{|c|}{\textbf{Unit}}\\ \hline
SSH &Dynamic sea surface height anomaly &m  \\ \hline
SSHIBC &The inverted barometer (ib) correction to sea surface height due to atmospheric pressure loading &m  \\ \hline
SSHNOIBC &Sea surface height anomaly without the inverted barometer (ib) correction &m  \\ \hline
\end{longtable}

\newp
\pagebreak
\subsection{dataset: SEA\_SURFACE\_HEIGHT, variable: SSH}
\begin{longtable}{|m{0.06\textwidth}|m{0.3\textwidth}|m{0.45\textwidth}|m{0.12\textwidth}|}
\caption{Attributes description for 'SSH'}
\label{tab:table-SEA_SURFACE_HEIGHT_SSH} \\ 
\hline \endhead \hline \endfoot
\rowcolor{lightgray} \textbf{Storage Type} & \textbf{Variable Name} & \textbf{Description} & \textbf{Unit} \\ \hline
float32 & SSH & Dynamic sea surface height anomaly & m \\ \hline
\multicolumn{4}{|c|}{\cellcolor{lightgray}{\textbf{Description of the variable in Common Data language (CDL)}}} \\ \hline
\multicolumn{4}{|c|}{\fontfamily{lmtt}\selectfont{\makecell{\parbox{.95\textwidth}{\vspace*{0.25cm} \footnotesize{float32 SSH(time, latitude, longitude)\\
\hspace*{0.5cm}SSH: \_FillValue = 9.96921e+36\\
\hspace*{0.5cm}SSH: coordinates = time\\
\hspace*{0.5cm}SSH: coverage\_content\_type = modelResult\\
\hspace*{0.5cm}SSH: long\_name = Dynamic sea surface height anomaly\\
\hspace*{0.5cm}SSH: standard\_name = sea surface height above geoid\\
\hspace*{0.5cm}SSH: units = m\\
\hspace*{0.5cm}SSH: valid\_max = 2.2875382900238037\\
\hspace*{0.5cm}SSH: valid\_min = -2.4861555099487305\\
}}}}} \\ \hline
\rowcolor{lightgray} \multicolumn{4}{|c|}{\textbf{Comments}} \\ \hline
\multicolumn{4}{|p{1\textwidth}|}{\footnotesize{{Dynamic sea surface height anomaly above the geoid, suitable for comparisons with altimetry sea surface height data products that apply the inverse barometer (ib) correction. note: ssh is calculated by correcting model sea level anomaly etan for three effects: a) global mean steric sea level changes related to density changes in the boussinesq volume-conserving model (greatbatch correction, see stergloh), b) the inverted barometer (ib) effect (see sshibc) and c) sea level displacement due to sea-ice and snow pressure loading (see siceload). ssh can be compared with the similarly-named ssh variable in previous ecco products that did not include atmospheric pressure loading (e.g., version 4 release 3). use sshnoibc for comparisons with altimetry data products that do not apply the ib correction.}}} \\ \hline
\end{longtable}

\begin{figure}[H]
\centering
\includegraphics[scale=0.55]{/Users/brucel/ecco/yip/ECCO-Dataset-Production/document_generator/general/files_general/version_specific/v4r4/output_and_granules/output/figures/figures_variables_latlon/Sea_Surface_Height/SSH.png}
\caption{dataset: SEA\_SURFACE\_HEIGHT, variable: SSH}
\label{tab:table-SEA_SURFACE_HEIGHT_SSH-Plot}
\end{figure}
\newpage
\pagebreak
\subsection{dataset: SEA\_SURFACE\_HEIGHT, variable: SSHIBC}
\begin{longtable}{|m{0.06\textwidth}|m{0.3\textwidth}|m{0.45\textwidth}|m{0.12\textwidth}|}
\caption{Attributes description for 'SSHIBC'}
\label{tab:table-SEA_SURFACE_HEIGHT_SSHIBC} \\ 
\hline \endhead \hline \endfoot
\rowcolor{lightgray} \textbf{Storage Type} & \textbf{Variable Name} & \textbf{Description} & \textbf{Unit} \\ \hline
float32 & SSHIBC & The inverted barometer (ib) correction to sea surface height due to atmospheric pressure loading & m \\ \hline
\multicolumn{4}{|c|}{\cellcolor{lightgray}{\textbf{Description of the variable in Common Data language (CDL)}}} \\ \hline
\multicolumn{4}{|c|}{\fontfamily{lmtt}\selectfont{\makecell{\parbox{.95\textwidth}{\vspace*{0.25cm} \footnotesize{float32 SSHIBC(time, latitude, longitude)\\
\hspace*{0.5cm}SSHIBC: \_FillValue = 9.96921e+36\\
\hspace*{0.5cm}SSHIBC: coordinates = time\\
\hspace*{0.5cm}SSHIBC: coverage\_content\_type = modelResult\\
\hspace*{0.5cm}SSHIBC: long\_name = The inverted barometer (IB) correction to sea surface height due to atmospheric pressure loading\\
\hspace*{0.5cm}SSHIBC: units = m\\
\hspace*{0.5cm}SSHIBC: valid\_max = 0.8955588340759277\\
\hspace*{0.5cm}SSHIBC: valid\_min = -0.5228679180145264\\
}}}}} \\ \hline
\rowcolor{lightgray} \multicolumn{4}{|c|}{\textbf{Comments}} \\ \hline
\multicolumn{4}{|p{1\textwidth}|}{\footnotesize{{Not an ssh itself, but a correction to model sea level anomaly (etan) required to account for the static part of sea surface displacement by atmosphere pressure loading: ssh = sshnoibc - sshibc. note: use ssh for model-data comparisons with altimetry data products that do apply the ib correction and sshnoibc for comparisons with altimetry data products that do not apply the ib correction.}}} \\ \hline
\end{longtable}

\begin{figure}[H]
\centering
\includegraphics[scale=0.55]{/Users/brucel/ecco/yip/ECCO-Dataset-Production/document_generator/general/files_general/version_specific/v4r4/output_and_granules/output/figures/figures_variables_latlon/Sea_Surface_Height/SSHIBC.png}
\caption{dataset: SEA\_SURFACE\_HEIGHT, variable: SSHIBC}
\label{tab:table-SEA_SURFACE_HEIGHT_SSHIBC-Plot}
\end{figure}
\newpage
\pagebreak
\subsection{dataset: SEA\_SURFACE\_HEIGHT, variable: SSHNOIBC}
\begin{longtable}{|m{0.06\textwidth}|m{0.3\textwidth}|m{0.45\textwidth}|m{0.12\textwidth}|}
\caption{Attributes description for 'SSHNOIBC'}
\label{tab:table-SEA_SURFACE_HEIGHT_SSHNOIBC} \\ 
\hline \endhead \hline \endfoot
\rowcolor{lightgray} \textbf{Storage Type} & \textbf{Variable Name} & \textbf{Description} & \textbf{Unit} \\ \hline
float32 & SSHNOIBC & Sea surface height anomaly without the inverted barometer (ib) correction & m \\ \hline
\multicolumn{4}{|c|}{\cellcolor{lightgray}{\textbf{Description of the variable in Common Data language (CDL)}}} \\ \hline
\multicolumn{4}{|c|}{\fontfamily{lmtt}\selectfont{\makecell{\parbox{.95\textwidth}{\vspace*{0.25cm} \footnotesize{float32 SSHNOIBC(time, latitude, longitude)\\
\hspace*{0.5cm}SSHNOIBC: \_FillValue = 9.96921e+36\\
\hspace*{0.5cm}SSHNOIBC: coordinates = time\\
\hspace*{0.5cm}SSHNOIBC: coverage\_content\_type = modelResult\\
\hspace*{0.5cm}SSHNOIBC: long\_name = Sea surface height anomaly without the inverted barometer (IB) correction\\
\hspace*{0.5cm}SSHNOIBC: units = m\\
\hspace*{0.5cm}SSHNOIBC: valid\_max = 2.2390522956848145\\
\hspace*{0.5cm}SSHNOIBC: valid\_min = -2.45104718208313\\
}}}}} \\ \hline
\rowcolor{lightgray} \multicolumn{4}{|c|}{\textbf{Comments}} \\ \hline
\multicolumn{4}{|p{1\textwidth}|}{\footnotesize{{Sea surface height anomaly above the geoid without the inverse barometer (ib) correction, suitable for comparisons with altimetry sea surface height data products that do not apply the inverse barometer (ib) correction. note: sshnoibc is calculated by correcting model sea level anomaly etan for two effects: a) global mean steric sea level changes related to density changes in the boussinesq volume-conserving model (greatbatch correction, see stergloh), b) sea level displacement due to sea-ice and snow pressure loading (see siceload). in ecco version 4 release 4 the model is forced with atmospheric pressure loading. sshnoibc does not correct for the static part of the effect of atmosphere pressure loading on sea surface height (the so-called inverse barometer (ib) correction). use ssh for comparisons with altimetry data products that do apply the ib correction.}}} \\ \hline
\end{longtable}

\begin{figure}[H]
\centering
\includegraphics[scale=0.55]{/Users/brucel/ecco/yip/ECCO-Dataset-Production/document_generator/general/files_general/version_specific/v4r4/output_and_granules/output/figures/figures_variables_latlon/Sea_Surface_Height/SSHNOIBC.png}
\caption{dataset: SEA\_SURFACE\_HEIGHT, variable: SSHNOIBC}
\label{tab:table-SEA_SURFACE_HEIGHT_SSHNOIBC-Plot}
\end{figure}
\newpage